\section{DGLAP evolution structure}

\subsection{When DGLAP is diagonal in pure gluodynamics}

The $M_{03} ^+\equiv
{ M}_{0 i; i 3} + { M}_{3 i;i 0}$ combination defined by Eq. (\ref{M03}) contains the twist-2 amplitude $\mathcal{M}_{pp}$,
\begin{align}
M_{03} ^+
= 4 p_0 p_3 \mathcal{M}_{pp} +
2 p_0 z_3 \left ( \mathcal{M}_{pz} +\mathcal{M}_{zp} \right) \ ,
\end{align}
though with a higher-twist admixture \mbox{$\mathcal{M}_{zp}+ \mathcal{M}_{pz}$}.
In the local limit, the relevant operator is proportional to the 3rd component of the Poynting vector
\begin{align*}
{\bf S_3}=
({\bf E \times B})_3 =E_1 B_2- B_1 E_2 =-( G_{01} G_{13} +G_{32} G_{20}) \ .
\end{align*}

As already mentioned, the box part of the one-loop correction to the matrix element
$M_{03} ^+(z_3,p)$
in pure gluodynamics
has a simple DGLAP structure\footnote{This simplicity may be violated in higher orders.} (\ref{box+}).
Combining all the gluon one-loop corrections to it,
we get, in the $\overline{\rm MS}$ scheme,
\begin{align}
& \frac{g^2N_c}{8\pi^2}\int_0^1 \dd u
\left\{ \left[ \left (\frac{3}{2} -\frac16 \right ) \ln(z_3^2 \mu_{\rm UV}^2 e^{2\gamma_E}/4) +2 \right] \delta(\bar u) \right. \nonumber \\
&\quad \left. - 2 \log(z_3^2 \mu_{\rm IR}^2 e^{2\gamma_E}/4) \left [ \frac{(1- u \bar u)^2}{\bar u} \right ]_+ \right.
\label{M031L} \\
&\quad \left. + \left [ u -3 \frac{u}{\bar u} -4 \frac{\log (\bar u)}{\bar u} \right ]_+
+ 2 \left (\bar uu + \frac{2}{3}\bar u^3 \right ) \right\} \ M_{03} ^+(uz,p) \ .\nonumber
\end{align}

The first line here comes from the UV-singular contributions. It contains the $\delta(\bar u)$ factor
which reflects
the local nature of the UV divergences and converts
$M_{03} ^+(uz,p)$ into $M_{03} ^+(z,p)$.
The second line contains the Altarelli-Parisi (AP) kernel
\begin{align}
B_{gg} (u) =&
2 \left [\frac{(1-u\bar u)^2} {1-u} \right ]_{+} \ .
\label{V1}
\end{align}
It has the plus-prescription
structure reflecting the fact that,
in the local limit, ${\mathcal{M}}_{pp}(z,p)$ is proportional to the
matrix element of the gluon energy-momentum tensor.
From now on, ``+'' means the plus-prescription at 1.

The third line
contains $z_3$-independent terms coming from the vertex diagrams
(these have the plus-prescription form) and from the box diagram.
The latter may be written as a sum of the term $2 \left (\bar uu + \frac{2}{3}\bar u^3 \right )_{+}$
that has the plus-prescription form and the term $\frac23 \delta (\bar u)$ that
may be combined with the UV terms.

\subsection{Reduced Ioffe-time distribution}

The combination $\mathcal{M}_{pz} +\mathcal{M}_{zp}$ is an odd function of $\nu =z_3 p_3$.
Writing it as $2 z_3 p_3 m^+_{zp} (\nu, z_3^2)$, we have
\begin{align}
M_{03} ^+(z_3, p) = 4 p_3 p_0 [ \mathcal{M}_{pp} (\nu, z_3^2 ) +
z_3^2 \ m^+_{zp} (\nu, z_3^2 )] \ .
\label{M03p}
\end{align}
Dividing out the kinematical factor $4 p_3 p_0$, we deal with
\begin{align}
\widetilde{\mathcal{M}}_{pp}(\nu, z_3^2 ) \equiv
\mathcal{M}_{pp}(\nu, z_3^2 ) +
z_3^2 \ m^+_{zp} (\nu, z_3^2 ) \ ,
\end{align}
which is a function of $\nu$ and $z_3^2$.
Now, just like in the quark case considered in Refs. \cite{Radyushkin:2017cyf,Orginos:2017kos}, we can introduce
the reduced Ioffe-time distribution
\begin{align}
\widetilde{\mathfrak M} (\nu, z_3^2) \equiv \frac{\widetilde{\mathcal{M}}_{pp}(\nu, z_3^2)}{\widetilde{\mathcal{M}}_{pp} (0, z_3^2)} \ .
\label{redm0}
\end{align}
Since $\widetilde{\mathcal{M}}_{pp}(\nu, z_3^2)$ is obtained from the multiplicatively
renormalizable combination $M_{03} ^+$,
the UV divergent $Z(z_3^2 \mu^2_{UV})$ factors generated by the link-related and gluon self-energy diagrams
cancel in the ratio (\ref{redm0}). As a result, the small-$z_3^2$ dependence of the reduced
pseudo-ITD $\widetilde{\mathfrak M} (\nu, z_3^2)$ comes from the logarithmic DGLAP evolution
effects only.
Moreover, $\widetilde{\mathcal{M}}_{pp}(0, z_3^2)$
has no DGLAP logarithmic dependence on $z_3^2$,
because of the plus-prescription nature of the AP kernel $B_{gg} (u)$.

Thus, neglecting ${\cal O} (z_3^2)$ terms, we conclude that,
{\it in pure gluodynamics, }
$\widetilde{\mathfrak M} (\nu, z_3^2)$ satisfies the evolution equation
\begin{align}
\frac{d}{d \ln z_3^2} \,
\widetilde{\mathfrak M} (\nu, z_3^2) &= - \frac{\alpha_s}{2\pi} \, N_c
\int_0^1 du \, B_{gg} (u) \widetilde{\mathfrak M} (u \nu, z_3^2)
\label{EE}
\end{align}
with respect to
$z_3^2$.
This relation is modified when gluon-quark transitions are present.

\subsection{Gluon-quark mixing}

In the $\overline{\rm MS}$ scheme, the contribution to $M_{03} ^+$ from the gluon-quark diagram shown in Fig. \ref{gluqum}
is given by
\begin{align}
& \frac{g^2 C_F}{4\pi^2 z_3} \int_0^1 \dd u
\Biggl [ -2u - \ln(z_3^2 \mu_{\rm IR}^2 e^{2\gamma_E}/4) \left[2\bar u+\delta(\bar u) \right] \Biggr ] {\cal O}_q (uz_3) \ ,
\label{gqmix}
\end{align}
where ${\cal O}_q (z_3)$ is a singlet combination of quark fields,
\begin{align}
& {\cal O}_q (z_3)= \frac{i}{2} \sum_f \left(\bar\psi_f( 0) \gamma_0 \psi_f (z_3) - \bar\psi_f( z_3) \gamma_0 \psi_f(0) \right) \ ,
\label{gsing}
\end{align}
with $f$ numerating quark flavors.
Note that ${\cal O}_q (z_3)$ vanishes
for $z_3=0$.
Expanding ${\cal O}_q (z_3)$ in $z_3$
\begin{align}
& {\cal O}_q ( z_3)= z_3\, \frac{i}{2} \sum_f \bar\psi_f( 0)\gamma_0 \!\stackrel{\leftrightarrow}{\partial}_3\! \psi_f(0) + {\cal O} (z_3^3) \ ,
\label{gsingb}
\end{align}
we see that the lowest term is proportional to the quark part of the energy-momentum tensor.
This term is accompanied by the $z_3$ factor which cancels the overall $1/z_3$ factor in Eq. (\ref{gqmix}).

The matrix element of ${\cal O}_q ( z_3)$ can be parametrized by
\begin{align}
& \langle p| {\cal O}_q ( z_3) |p\rangle = 2 p^0
\int_0^1 \dd x \, \sin (x p_3 z_3) \, q_S (x)
\label{qpar}
\end{align}
where $f_S (x) = \sum_f [q_f (x) +\bar q_f (x)]$ is the singlet quark distribution.
To extract the overall $z_3$ factor, we rewrite
\begin{align}
& \langle p| {\cal O}_q ( z_3) |p\rangle
= 2 p^0 p_3 z_3
\int_0^1 \dd y \, \int_0^1 \dd \alpha\, \cos (\alpha y \nu) \, y \, f_S (y )
\label{qpar2}
\end{align}
where $\nu = p_3 z_3$, as usual. This gives
\begin{align}
&\frac1{z_3} \int_0^1 \dd u \, A(u) \, \langle p| {\cal O}_q (u z_3) |p\rangle
\nn &
= 2 p^0 p_3 \, \int_0^1 \dd w {\cal I}_S (w \nu) {\cal A} (w)
\label{qpar3}
\end{align}
for $u$-integrals of Eq. (\ref{gqmix}) type. Here
\begin{align}
& {\cal I}_S (\nu) =
\int_0^1 \dd y \ \cos (y \nu) \, y\, f_S (y)
\label{MS}
\end{align}
is the singlet quark Ioffe-time distribution, and
\begin{align}
{\cal A} (w)=
\int_w^1 \dd u \, A(u) \ .
\label{Aw}
\end{align}
For the evolution kernel $B_{gq} (u) \equiv 2\bar u+\delta(\bar u)$, we get
\begin{align}
{\cal B}_{gq} (w)=
\int_w^1 \dd u \, B_{gq} (u) = 1+ (1- w)^2 \ .
\label{Bw}
\end{align}

\subsection{Matching relations}

A disadvantage of the $M_{03} ^+(z_3, p)$ combination is that
it vanishes
when $p_3=0$ (see \mbox{Eq. (\ref{M03p}))}.
Thus, to extract
$\widetilde{\mathcal{M}}_{pp}(\nu, z_3^2 )$
for $\nu=0$,
one should
make measurements of $M_{03} ^+(z_3, p)$ for a few low values of $p_3$,
divide $p_3$ out,
and extrapolate the results to \mbox{$p_3=0$.} This procedure leads
to additional systematic uncertainties.

\begin{figure}[t]
\centerline{\includegraphics[width=1.7in]{images/gluqum}}
\vspace{-3mm}
\caption{Gluon-quark mixing diagram. \label{gluqum}}
\end{figure}

Fortunately,
the combination
${ M}_{0 i; i 0} - { M}_{i j; i j} = 2 p_0^2 \mathcal{M}_{pp}$ of \mbox{Eq. (\ref{00m}),} being proportional
to $p_0^2$, does not have this problem.
Furthermore, it gives the twist-2 amplitude $\mathcal{M}_{pp}$ without contaminations.
The amplitude $\mathcal{M}_{pp} (\nu, z_3^2)$ obtained in this way
may be used
to form the reduced pseudo-ITD
${\mathfrak M} (\nu, z_3^2)$, as in Eq. (\ref{redm0}).

Using the results of our calculations for the one-loop corrections to
${ M}_{0 i; i 0}$ and ${ M}_{i j; i j}$, and keeping just the ${\cal M}_{pp}$ term in the correction (while skipping the ``higher twist'' terms
${\cal M}_{zz}, {\cal M}_{zp}, {\cal M}_{pz}, {\cal M}_{ppzz}$)
we obtain
the matching relation
\begin{align}
{\mathfrak M}& (\nu, z_3^2)\, {{\cal I}_g (0, \mu^2) } = { {\cal I}_g (\nu, \mu^2) }
- \frac{\alpha_s N_c}{2\pi}\int_0^1 \dd u \, {{\cal I}_g (u \nu, \mu^2)}
\nonumber \\
& \times
\Biggl \{ \ln (z_3^2 \mu^2 e^{2\gamma_E}/4) \, B_{gg} (u) +4 \left [ \frac{u+ \log (\bar u)}{\bar u}\right ]_{+} \nonumber \\
& + \frac23 \left [ 1 -u^3
\right ]_+ \Biggr \}
- \frac{\alpha_s C_F}{2\pi} \ln (z_3^2 \mu^2 e^{2\gamma_E}/4)
%
\nonumber \\
&
\times
\int_0^1 \dd w \,
\left [{\cal I}_S (w \nu, \mu^2) - {\cal I}_S (0, \mu^2) \right ]
\, {\cal B}_{gq} (w) \
\label{matching}
\end{align}
between the ``lattice function''
${\mathfrak M}( \nu,z_3^2)$ and the light-cone ITDs ${\cal I}_g (\nu,\mu^2)$ and
${\cal I}_S(\nu,\mu^2)$. The first of them is
related to the gluon PDF ${f}_g(x,\mu^2)$
by
\begin{align}
{\cal I}_g (\nu,\mu^2) =\frac12
\int_{-1}^1 dx \, \,
e^{ix\nu} \,x {f}_g(x,\mu^2) \ .
\label{If}
\end{align}
Since $x {f}_g(x,\mu^2)$ is an even function of $x$,
the real part of ${\cal I}_g (\nu,\mu^2)$ is given by the cosine transform of $x {f}_g(x,\mu^2)$,
while its imaginary part vanishes.
The factor ${\cal I}_g (0, \mu^2)$ has the meaning of the fraction of the hadron
momentum carried by the gluons, ${\cal I}_g (0, \mu^2) = \langle x \rangle_{\mu^2}$.

Thus, Eq. (\ref{matching}) allows to extract the shape of the gluon distribution.
Its normalization, i.e., the value of $\langle x \rangle_{\mu^2}$ should be
found by an independent lattice calculation, similar to that performed
in Ref. \cite{Yang:2018bft}.
The singlet quark function ${\cal I}_S (w \nu, \mu^2)$ that appears in the
${\cal O} (\alpha_s)$ correction should be also calculated (or estimated)
independently.

Substituting Eq. (\ref{If}) into the matching condition (\ref{matching}), we can
rewrite the latter in the kernel form \cite{Radyushkin:2019owq}
\begin{align}
{\mathfrak M}(\nu,z_3^2) & =
\int_{0}^1 \dd x \,
\frac{x {f}_g (x,\mu^2) }{\langle x \rangle_{\mu^2}} \, R_{gg} (x \nu, z_3^2 \mu^2 ) \nn &
+
\int_{0}^1 \dd x \,
\frac{x {f}_S (x,\mu^2) }{\langle x \rangle_{\mu^2}} \, R_{gq} (x \nu, z_3^2 \mu^2 ) \
,
\label{MtchI}
\end{align}
where the kernel $R_{gg} (x \nu, z_3^2 \mu^2 )$
is given by
\begin{align}
R_{gg}(y, z_3^2 \mu^2) =& \cos y
- \frac{\alpha_s}{2\pi} \, N_c
\Biggl \{ \ln \left ( z_3^2 \mu^2\frac{e^{2\gamma_E+2}}{4} \right ) \,
R_B(y)\nn & +R_U (y)
+R_L(y) +R_C(y) \Biggr \} \, ,
\label{MtchR}
\end{align}
with $R_B(y)$ being
the cosine Fourier transform of the $B_{gg}$ kernel
\begin{align}
R_B(y) =& \int_0^1 \dd u \, B_{gg} (u) \, \cos (u y) \ .
\label{RB}
\end{align}
Its calculation is straightforward, and the result is expressed in terms of $\cos y$, $\sin y$
and the integral cosine $\text{Ci}(y)$ and sine $\text{Si}(y)$ functions.
The latter come from the $1/(1-u)$ part of $B(u)$, which gives
\begin{align}
\int_0^1 \dd u \, & \left[ \frac1{1-u} \right ]_+ \, \cos (u y) = \sin (y) \, \text{Si}(y) \nn & +\cos (y)\,
[\text{Ci}(y)-\log (y)-\gamma_E ]
\ .
\label{RB0}
\end{align}
This combination also appears in $R_U(y)$, which is given by the cosine transform of $4[u/\bar u]_+$.
Similarly, $R_L(y)$ is the cosine transform of the \mbox{$4[(\ln \bar u)/\bar u]_+$} term.
It involves a hypergeometric function
\begin{align}
R_L(y) =
4\, {\rm Re} \left [ i y e^{i y} \,
_3F_3(1,1,1;2,2,2;-i y) \right ] \, .
\label{MtchRL}
\end{align}

The $R_C(y)$ and $R_{gq} (y)$ kernels are given by the cosine transforms of $\frac23 \left [ 1-u^3
\right ]_+$ and $\left [ 1+\bar u^2\right]_+$, respectively.
Expressions for them
involve $\cos y$, $\sin y$ and inverse powers of $y$.

The important point is that the $R(y, z_3^2 \mu^2)$ kernels are given by {\it explicit} perturbatively calculable
expressions. Using them and Eq. (\ref{MtchI}) one may directly relate
${\mathfrak M}(\nu,z_3^2)$ and the light-cone
PDFs.
Then,
assuming some parameterizations for the $f_g(x,\mu^2)$ and $f_S(x, \mu^2)$ distributions,
one can fit their parameters and $\alpha_s$
from the lattice data for ${\mathfrak M}(\nu,z_3^2)$ using Eqs. (\ref{MtchI}), (\ref{MtchR}). This procedure is essentially the same as that used
in the ``good lattice cross sections'' approach \cite{Ma:2014jla,Ma:2017pxb}.

\subsection{Matching relations for quasi-PDFs}

The kernel relation (\ref{MtchI}) directly connects ${\mathfrak M}(\nu,z_3^2)$
with PDFs.
So, there is no need to introduce
intermediate functions, such as quasi-PDFs.
Still, our results for particular matrix elements,
such as Eq. (\ref{M031L}) for $M_{03} ^+(z_3,p)$,
may be used to get matching conditions for quasi-PDFs. The latter
are generically defined \cite{Ji:2013dva} as
\begin{align}
Q(y, p_3) = \frac{p_3}{2 \pi}
\int_{-\infty}^\infty \dd z_3
{\cal M} (z_3, p) \, e^{-i y p_3 z_3} \
.
\label{qPDF}
\end{align}

To proceed, one should write the amplitudes ${\cal M} (z_3, p)$ through the kernel
relation (\ref{MtchI}) with $R(x \nu, z_3^2 \mu^2 )$ expressed in terms of $p_3$ and $z_3$,
call it $J(x, p_3, z_3)$.
The structure of its dependence on $z_3$ at one loop may be read off Eq. (\ref{M031L}).
For the $gg$ part,
\begin{align}
& J_1^{gg} (x, p_3, z_3)= \left (\gamma_U \ln z_3^2 +C_U \right) e^{ix p_3 z_3}
\nn & - \int_0^1 \dd u \,\left [ \ln z_3^2 \, B_{gg} (u)
+C (u) \right ]e^{iu x p_3 z_3} \ .
\label{Jk}
\end{align}
The 1-loop quasi-PDF matching kernel is then given by
\begin{align}
Z_1^{gg} (y, x, p_3) = \frac{p_3}{2 \pi}
\int_{-\infty}^\infty \dd z_3 \, J_1^{gg} (x, p_3, z_3)
\, e^{-i y p_3 z_3} \
.
\label{qZ}
\end{align}
The $C_U$ and $C(u)$ contributions of $J_1^{gg} (x, p_3, z_3)$
produce \mbox{$C_U \delta (y- x)$} and $C (u)\delta (y - u x )$ terms.
Hence,
the resulting parts of $Q(y,p_3)$ are visible
in the ``canonical'' $0\leq |y| \leq 1$ region only. However, the terms
with $\ln z_3^2$ give nonzero contributions in the $y/x>1$
and $y/x<0$ regions
as well, namely
\begin{align}
& \, Z_1^{gg} (y, x, p_3) |_{y/x>1} = \frac1{|x|} \left [- \frac{\gamma_U}{\eta -1}
+ \int_0^1 \dd u \, \frac{B_{gg} (u)}{\eta-u} \right ]
\ ,
\label{Zy}
\end{align}
where $\eta= y/x$.
The $Z_1^{gg} (y, x, p_3) |_{y/x<0}$ term is given by the same expression, but with the opposite sign.
Note that these contributions are completely
determined by the AP kernel $B_{gg} (u)$ and the UV constant
$\gamma_U$.
Knowing them, one derives from Eq. (\ref{Zy}) a general constraint on the results
for $Z_1^{gg} (\eta, 1, p_3) |_{\eta>1}$
and $Z_1^{gg} (\eta, 1, p_3) |_{\eta<0}$
obtained by any Feynman diagram calculation.
Using explicit form of $B_{gg} (u)$, we find
\begin{align}
\int_0^1 \dd u \, \frac{B_{gg} (u) }{\eta-u} & = 2 \frac{(1-\eta \bar \eta)^2}{\eta -1} \, \ln \frac{\eta-1}{\eta}
\nn & + \frac{11}{6} \frac{1}{\eta -1} + \eta (2 \eta -1) +\frac{11}{3}
\ .
\label{ZB}
\end{align}

For large $\eta$, this expression tends to zero as ${\cal O}(1/\eta^2)$.
It should be stressed that such a behavior results from any kernel $B(u)$ that has the plus-prescription form.

This observation and the explicit expression given by \mbox{Eq. (\ref{ZB})}
may be used to check the gluon-gluon matching kernels
in \mbox{Refs. \cite{Wang:2017qyg,Wang:2019tgg}. }
Our check shows that $Z_1^{gg} (y, 1, p_3) |_{y>1}$ corresponding to Eq. (64) of Ref.
\cite{Wang:2019tgg} does not satisfy the constraint (\ref{Zy}).
The difference is by a constant term (-2/3) that leads to a linear divergence
in the integral of $Z_1^{gg} (y, 1, p_3) |_{y>1}$ \mbox{over $y$. }
The same difference (with the opposite sign) appears
in the $Z_1^{gg} (y, 1, p_3) |_{y<0}$ term in Eq. (64) of Ref. \cite{Wang:2019tgg}.
Apparently, these differences result from the use of
the off-shell external gluon fields in the calculations of Ref. \cite{Wang:2019tgg},
but a discussion of this topic is outside of the scope of our paper.
